% Generated draft: requires mathematical and source review.
\documentclass[11pt]{article}
\usepackage[margin=25mm]{geometry}
\usepackage{fontspec}
\setmainfont{Latin Modern Roman}
\usepackage{amsmath,amssymb,mathtools,unicode-math}
\setmathfont{Latin Modern Math}
\usepackage{enumitem,needspace,xurl,hyperref}
\setlength{\parindent}{0pt}
\setlength{\parskip}{5pt}
\setlength{\emergencystretch}{3em}
\newcommand{\N}{\mathbb{N}}
\newcommand{\Z}{\mathbb{Z}}
\newcommand{\Q}{\mathbb{Q}}
\newcommand{\R}{\mathbb{R}}
\newcommand{\C}{\mathbb{C}}
\newcommand{\F}{\mathbb{F}}
\newcommand{\A}{\mathbb{A}}
\newcommand{\PP}{\mathbb{P}}
\newcommand{\E}{\mathbb{E}}
\newcommand{\eps}{\varepsilon}
\newcommand{\dd}{\,\mathrm{d}}
\begin{document}
\section*{UnsolvedMath ID: 30005185}
\textbf{Source problem code:} catalogue entry
\subsection{Definitions and mathematical statement}
Let \(p\) be a prime and let \(R\in\{\mathbb F_p,\mathbb Z\}\). Use the conventional smooth interpretation: knots are smooth embeddings of \(S^1\) in \(S^3\), considered up to ambient isotopy. Two knots \(K_0,K_1\) are smoothly concordant if they cobound a smoothly and properly embedded annulus in \(S^3\times[0,1]\), with boundary knots in the respective boundary components. A knot is smoothly slice if it bounds a smoothly and properly embedded disk in \(B^4\). A ribbon knot bounds such a disk admitting a ribbon presentation, with only index-zero and index-one critical points for the radial function in its interior. The supplied statement does not explicitly specify the concordance category; the smooth convention therefore requires confirmation against the cited report.

Write \(\widetilde{\operatorname{Kh}}^{i,j}(K;R)\) for reduced Khovanov homology with coefficients in \(R\), where \(i\) and \(j\) are the homological and quantum gradings. Interpret its total rank by
\[
 r_R(K)=\sum_{i,j}\operatorname{rank}_R\widetilde{\operatorname{Kh}}^{i,j}(K;R),
\]
where rank over \(\mathbb F_p\) means vector-space dimension and
\[
 \operatorname{rank}_{\mathbb Z}M=\dim_{\mathbb Q}(M\otimes_{\mathbb Z}\mathbb Q).
\]
Thus the integral interpretation counts free rank and does not count torsion generators. These are conventional interpretations of the unqualified rank notation in the supplied record.

\textbf{Proposed concordance-invariance statement.} For every coefficient choice above and every pair of smoothly concordant knots \(K_0,K_1\), prove that
\[
 r_R(K_0)\equiv r_R(K_1)\pmod 4.
\]
Equivalently, the proposed function \(K\mapsto r_R(K)\bmod 4\) should descend to smooth knot-concordance classes.

\textbf{Ribbon-knot formulation supplied by the registry.} For every coefficient choice above and every ribbon knot \(J\), prove that
\[
 r_R(J)\equiv 1\pmod 4.
\]
The registry explicitly presents this ribbon-knot assertion as equivalent to the concordance-invariance statement. That equivalence is retained as a source assertion; the supplied excerpt does not contain its justification or the hypotheses needed to check it.

The original passage motivates a modulo-four weakening of assertions about slice knots, referring to Conjectures 2.(2) and 1.(2). With the usual normalization \(r_R(U)=1\) for the unknot \(U\), the concordance-invariance proposal would in particular imply \(r_R(K)\equiv1\pmod4\) for every smoothly slice knot \(K\). The referenced conjectures are not reproduced in the supplied record, so their full scope cannot be reconstructed here.
\subsection{Short English statement}
The proposed conjecture says that the total rank of reduced Khovanov homology, modulo four, depends only on the smooth concordance class of a knot, for coefficients in every prime field or in \(\mathbb Z\). The registry also supplies a ribbon-knot congruence and asserts its equivalence to this proposal. The smooth category, integral-rank convention, and claimed equivalence require source review. This is a statement-only draft.
\subsection{Sources}
\begin{itemize}
\item Supplied local UnsolvedMath registry record, permanent ID \(30005185\), title \emph{Modulo-Four Reduced Khovanov Rank of Ribbon Knots}, source problem code \texttt{catalogue entry}, in the \emph{Oberwolfach Reports Open Problems} collection. No individual proposer is named; the supplied proposed year is 2022.
\item Supplied source citation: \emph{Topologie (2023)}, Oberwolfach Reports, DOI 10.4171/owr/2022/34, pp. 1927--1986. Original Oberwolfach report: \url{https://doi.org/10.4171/owr/2022/34}. The supplied background describes the reports as records of mathematical workshops and their open-problem discussions. The citation's 2023 label and the record's 2022 proposed year are retained as supplied.
\item Supplied literature source attributed to Dunfield--Gong--Hockenhull--Marengon--Willis: \url{https://arxiv.org/abs/2303.04233}. The record supplies no paper title or publication details beyond this attribution and link.
\item Literature assessment dated exactly \texttt{2026-08-22}: the supplied curation labels the modulo-four proposal for slice or ribbon knots open and conjectural. It reports counterexamples to the modulo-eight version in the follow-up paper, with evidence and special cases for the modulo-four replacement. This is supplied dated curation, not an independently verified literature assessment; its status and cited claims require independent verification.
\item The registry labels the statement \texttt{corrected\_verified} and states that it was reconstructed and checked against the official Oberwolfach report. These are supplied verification claims, not verification performed for this draft. The original excerpt refers to Conjectures 2.(2) and 1.(2), an odd-square condition, and a slice-knot Floer-rank example with reported rank 33, but contains damaged mathematical notation.
\end{itemize}
\end{document}
